\newpage
\section*{Đề bài}
\addcontentsline{toc}{section}{ĐỀ BÀI}

Dựa vào lý thuyết “Chapter 4 – Logistic Regression” đã học. Hãy tính toán và thực hiện các yêu cầu bên dưới như sau:

\textbf{Yêu cầu:}

Bạn có file dữ liệu Death\_Probability.csv đính kèm, trong đó có 03 cột: ID (bệnh nhân), PROCALCITONIN (máu, ng/ml), DEATH (0: sống; 1: tử vong).

\begin{enumerate}
\item Dựa vào công thức hồi quy logistic:
\[
p = \dfrac{1}{1 + e^{-(\beta_0 + \beta_1 x_1 + \beta_2 x_2 + \cdots + \beta_q x_q)}}
\]

Hãy ước lượng các tham số và đưa ra phương trình hồi quy logistic dựa vào tập dữ liệu đã cho.

\item Dựa vào công thức Odds:
\[
\textbf{Odds}(Y=1) = e^{\beta_0 + \beta_1 x_1 + \beta_2 x_2 + \cdots + \beta_q x_q}
\]
Hãy diễn giải, so sánh và tiên lượng khả năng tử vong của bệnh nhân, khi PROCALCITONIN = 0 ng/ml và PROCALCITONIN = 1 ng/ml.

\end{enumerate}

\textbf{Chú ý:}

Kết quả phần (A) phải được tóm tắt và chụp hình đưa vào tệp tin có quy cách đặt tên như sau: MSSV\_HọTên\_Lab04\_Output.docx (.pdf)

\begin{itemize}

\item Gửi đính kèm các file SPSS, R và Python của Phần A. Tương tự, quy cách đặt tên file như sau:
\begin{itemize}
\item MSSV\_HọTên\_Lab04\_SPSS.sav (và .spv)
\item MSSV\_HọTên\_Lab04\_Python.ipynb
\item MSSV\_HọTên\_Lab04\_R.r
\end{itemize}

\item Đóng gói lại thành 01 tập tin duy nhất và nộp file này theo đúng deadline với quy cách đặt tên file như sau: MSSV\_HọTên\_Lab04.rar (.zip)

\item Deadline nộp bài: 31/03/2026

\item Sinh viên nộp bài Lab04 không đúng deadline, xem như không có bài Lab04.

\item Sinh viên nộp không đủ file và đặt tên không đúng quy cách tên file, sẽ bị trừ điểm.

\item Các bài có hình thức trình bày và nội dung giống nhau sẽ nhận điểm Zero.

\end{itemize}